% !TEX root = ../main.tex
\begin{ouverture}[listes]{12}{Algorithmique et programmation}{Partie 5 · Méthodes transversales}
{Variables\enspace·\enspace Conditions\enspace·\enspace Boucles\enspace·\enspace Fonctions et listes}
\objectif{Lire et écrire des programmes Python courts.}
\objectif{Créer, parcourir et modifier une liste.}
\objectif{Programmer des termes, des sommes et des seuils.}
\objectif{Décomposer une simulation en fonctions et vérifier ses résultats.}
\prerequis{Calcul numérique et suites.}
\prerequis{Les exercices de simulation utilisent les chapitres 9 à 11.}
\end{ouverture}
\essentiel
\essSection{1}{Variables, tests et fonctions}
\begin{definition}
Une affectation comme \code{u = 3} associe une valeur à un nom. En langage
naturel, on peut la noter $u\leftarrow3$. Le signe \code{==} teste l'égalité.
Les types courants sont les entiers, les nombres décimaux et les booléens.
Une fonction définie avec \code{def} reçoit des arguments et renvoie une
valeur avec \code{return}. L'indentation organise les blocs.
\end{definition}
\begin{propriete}[Choisir une boucle]
La boucle \code{for k in range(n)} effectue $n$ passages, avec $k$ de 0 à $n-1$.
La boucle \code{while condition} continue tant que la condition est vraie.
Il faut actualiser les variables de la condition et étudier l'arrêt.
Un \code{if} exécute un bloc selon un test ; \code{else} traite l'autre cas.
\end{propriete}
\begin{savoirfaire}[Suivre l'état d'un programme]
Construire un tableau avec une ligne par passage et une colonne par variable.
Distinguer la valeur avant le passage et la valeur après l'affectation.
Pour un compteur, préciser ce qu'il compte : termes, itérations ou succès.
\end{savoirfaire}
\begin{center}
\begin{tikzpicture}[node distance=11mm]
\node[otBloc](a){Initialiser $u$ et $n$};
\node[otTest,below=of a](b){Seuil atteint ?};
\node[otBloc,right=24mm of b](c){Actualiser $u$ et $n$};
\node[otBorne,below=of b](d){Renvoyer $n$};
\draw[otFleche](a)--(b);\draw[otFleche](b)--node[right,font=\scriptsize]{oui}(d);
\draw[otFleche](b)--node[above,font=\scriptsize]{non}(c);
\coordinate (reprise) at ([yshift=4mm]b.north);
\draw[otFleche](c.north)|-(reprise);
\fill[otMuted](reprise)circle(.6pt);
\end{tikzpicture}
\par\legende{Une boucle de seuil demande une initialisation et une actualisation cohérentes.}
\end{center}

\clearpage
\essSection{2}{Créer et manipuler des listes}
\begin{definition}
Une liste regroupe des valeurs ordonnées : \code{L = [3, 1, 4]}.
Les indices commencent à 0 ; \code{L[0]} vaut 3 et \code{L[-1]} vaut 4.
\code{len(L)} donne la longueur et \code{sum(L)} la somme.
\end{definition}
\begin{propriete}
\code{L.append(x)} ajoute $x$ à la fin ; \code{L[i] = x} modifie une valeur ;
\code{L.pop(i)} supprime et renvoie la valeur d'indice $i$.
On peut parcourir directement les valeurs avec \code{for x in L}.
On peut aussi parcourir les indices avec \code{for i in range(len(L))}.
\end{propriete}
\begin{savoirfaire}[Compréhension et filtrage]
\code{[k*k for k in range(6)]} crée les carrés de 0 à 5.
\code{[x for x in L if x > 0]} sélectionne les éléments positifs.
Une liste peut servir à stocker les termes d'une suite, des valeurs d'une fonction
ou les résultats d'une simulation.
\end{savoirfaire}
\essSection{3}{Tester et organiser un programme}
\begin{propriete}
Décomposer une tâche en fonctions rend chaque étape contrôlable.
Tester les cas usuels et les cas limites autorisés : $n=0$, liste à un élément,
seuil déjà atteint. Vérifier les hypothèses : dénominateur non nul, liste non vide,
entier positif. Une simulation ne démontre pas une propriété générale.
\end{propriete}
\begin{attention}
Une boucle sur \code{range(n)} ne va pas jusqu'à $n$ inclus. Pour calculer
$u_0,\ldots,u_n$, il y a $n$ transitions mais $n+1$ valeurs à conserver.
Le signe \code{**} désigne une puissance ; \code{\textasciicircum} ne la remplace pas en Python.
\end{attention}

\exercices

\rubrique{Découverte}
\begin{blocLivre}{ex-12-1}
\begin{exercice}[1]{Suivre des affectations}
Exécuter mentalement le code, puis donner les valeurs finales de \code{a} et \code{b}.
\begin{python}
a = 3
b = 5
a = a + b
b = a - b
a = a - b
\end{python}
Quel effet global a ce programme ?
\end{exercice}
\end{blocLivre}
\begin{blocLivre}{ex-12-2}
\begin{exercice}[1]{Tester une condition}
Que renvoie la fonction pour 7, 8 et 0 ? Expliquer le test.
\begin{python}
def traitement(n):
    if n % 2 == 0:
        return n//2
    return 3*n + 1
\end{python}
\end{exercice}
\end{blocLivre}
\begin{blocLivre}{ex-12-3}
\begin{exercice}[1]{Définir une fonction}
Écrire une fonction Python pour $f(x)=x^2-3x+2$, puis calculer \code{f(0)}, \code{f(2)} et \code{f(5)}. Faut-il modifier le nom de la variable dans la formule mathématique ?
\end{exercice}
\end{blocLivre}
\begin{blocLivre}{ex-12-4}
\begin{exercice}[1]{Modifier une liste}
Donner l'état de la liste après chaque ligne et les valeurs finales de \code{len(L)} et \code{sum(L)}.
\begin{python}
L = [3, 1, 4]
L.append(2)
L[1] = 5
L.pop(0)
\end{python}
\end{exercice}
\end{blocLivre}
\begin{blocLivre}{ex-12-5}
\begin{exercice}[1]{Comprendre une liste en compréhension}
Donner les deux listes créées, puis leurs longueurs.
\begin{python}
L = [k*k for k in range(6)]
P = [x for x in L if x % 2 == 0]
\end{python}
\end{exercice}
\end{blocLivre}
\clearpage
\rubrique{Entraînement}
\begin{blocLivre}{ex-12-6}
\begin{exercice}[2]{Calculer une factorielle}
Pour un entier $n\geq0$, écrire une fonction renvoyant $n!$, avec $0!=1$.\begin{questions}
\item Quel produit accumuler à chaque passage ?
\item Calculer les résultats pour 0, 1 et 5.
\item Pourquoi une initialisation à 0 serait-elle fausse ?
\end{questions}
\end{exercice}
\end{blocLivre}
\begin{blocLivre}{ex-12-7}
\begin{exercice}[2]{Accumuler une somme}
Écrire une fonction \code{somme(n)} qui calcule $1+2+\cdots+n$ pour $n\geq0$, par une boucle. Vérifier pour $n=0$ et $n=10$, puis comparer à $n(n+1)/2$.
\end{exercice}
\end{blocLivre}
\begin{blocLivre}{ex-12-8}
\begin{exercice}[2]{Trouver un premier seuil}
Une suite vérifie $u_0=100$ et $u_{n+1}=u_n+30$.\begin{questions}
\item Écrire une boucle qui renvoie le premier rang tel que $u_n\geq500$.
\item Vérifier le résultat avec la formule explicite.
\item Que renverrait le même principe si le seuil était déjà atteint au rang 0 ?
\end{questions}
\end{exercice}
\end{blocLivre}
\begin{blocLivre}{ex-12-9}
\begin{exercice}[2]{Additionner les termes d'une suite}
On veut calculer les $n$ premiers termes de la suite $u_k=2\times3^k$, de $k=0$ à $k=n-1$.\begin{questions}
\item Écrire une fonction de somme par actualisations successives.
\item Calculer le résultat pour $n=4$.
\item Donner une formule et traiter le cas $n=0$.
\end{questions}
\end{exercice}
\end{blocLivre}
\begin{blocLivre}{ex-12-10}
\begin{exercice}[2]{Compter les lettres sans autre collection}
Pour un texte non vide, écrire une fonction qui renvoie la fréquence d’une lettre, en parcourant le texte avec une boucle. Tester sur \code{BANANE} pour A et B.
\end{exercice}
\end{blocLivre}
\begin{blocLivre}{ex-12-11}
\begin{exercice}[2]{Calculer les moments d'une loi}
Les listes \code{valeurs} et \code{probas} ont la même longueur et décrivent une loi de probabilité valide. Écrire une fonction renvoyant espérance et variance, puis la tester sur $[-1,0,2]$ et $[0{,}2,0{,}5,0{,}3]$.
\end{exercice}
\end{blocLivre}
\clearpage
\rubrique{Problèmes}
\begin{blocLivre}{ex-12-12}
\begin{exercice}[3]{Construire la suite de Fibonacci}
On définit $F_0=0$, $F_1=1$ et $F_{n+2}=F_{n+1}+F_n$.\begin{questions}
\item Écrire une fonction qui renvoie la liste de $F_0$ à $F_n$, pour $n\geq0$.
\item Donner cette liste pour $n=7$.
\item Chercher le premier rang où le terme dépasse 1000 et vérifier le terme précédent.
\end{questions}
\end{exercice}
\end{blocLivre}
\begin{blocLivre}{ex-12-13}
\begin{exercice}[3]{Expérimenter avec Syracuse}
À partir d’un entier positif, on le divise par 2 s’il est pair ; sinon on le remplace par $3n+1$.\begin{questions}
\item Donner les termes à partir de 6 jusqu’à 1.
\item Écrire une fonction qui s’arrête à 1 ou après un nombre maximal d’étapes.
\item Pourquoi la borne sur le nombre d’étapes garantit-elle la fin du programme, indépendamment d’une propriété générale de la suite ?
\end{questions}
\end{exercice}
\end{blocLivre}
\begin{blocLivre}{ex-12-14}
\begin{exercice}[3]{Une approximation par Newton}
On cherche une valeur approchée de $\sqrt2$, racine positive de $f(x)=x^2-2$.\begin{questions}
\item Pour $a>0$, montrer que la tangente à la courbe de $f$ au point d’abscisse $a$ coupe l’axe des abscisses au point d’abscisse $(a+2/a)/2$. La méthode de Newton pose donc $x_0=1$ et $x_{k+1}=(x_k+2/x_k)/2$.
\item Écrire une fonction calculant $x_n$.
\item Calculer $x_1$ et $x_2$ exactement.
\item Expérimenter avec $n=5$ et comparer le carré obtenu à 2.
\item Vérifier que les termes calculés restent positifs, ce qui autorise la division.
\end{questions}
\end{exercice}
\end{blocLivre}
\begin{blocLivre}{ex-12-15}
\begin{exercice}[3]{Un mini-projet Monte-Carlo}
On veut estimer simultanément l’aire sous $y=x^2$ dans le carré unité et le nombre $\pi$.\begin{questions}
\item Écrire une fonction qui tire $n\geq1$ points, avec une graine fixée, et renvoie les deux estimations.
\item Expliquer pourquoi on peut réutiliser les mêmes points pour les deux comptes.
\item Proposer des contrôles automatiques simples des valeurs renvoyées.
\end{questions}
\label{figlivre-r12-15-zones-ex}
\figureLivre{r12-15-zones}

\end{exercice}
\end{blocLivre}
